\scanpage{29}%
\mspage{14}\space
$2$ nilpotent, ce qui permet de mettre la courbe sous la forme (4.1)
\[
y^2+xy=x^3+a_2\,x+a_6
\]
et un isomorphisme d'une telle courbe sur une autre est décrit par $s$ vérifiant (4.3) pour $u=1+2s$. Lorsque $j=j'$, la première de ces équations
\[
(1+2s)^2\,a'_2=a_2-s-s^2
\]
définit un revêtement étale de degré deux, et elle implique la seconde ; pour le voir, on se ramène à supposer $a_2=a'_2$ et on remarque que d'après (4.4)
\[
1/j=a_6+2\,R(a_2,a_6)
\]
pour une fraction rationnelle $R$ .

Pour vérifier (II) dans les cas non couverts par (III), on peut supposer que $S$ est le spectre d'un corps algébriquement clos et que $c_4=0$ ou $c_6=0$ \tikz[baseline=(punct.base)]{\node[inner sep=0pt] (punct) {,};\node[inner sep=0pt,anchor=north] at ([yshift=-0.16em]punct.south) {.};} i.e. que $j=0$ ou $j=1728$. Soit $p$ la caractéristique du corps,

a) $p\ne 2,3$ : on met la courbe sous la forme (2.1) ; si $c_4=0$ (resp. $c_6=0$) on aura $a_4=0$ (resp. $a_6=0$) et on peut par homothétic transformer la courbe en la courbe type
\[
y^2=x^3+1 \qquad (\text{si }c_4=0) \tag*{(6.5)}
\]
\[
y^2=x^3-x \qquad (\text{si }c_6=0) \tag*{(6.6)}
\]

b) $\unsure{\mathbb{\text{\makebox[0pt][l]{\kern.45pt$p$}$p$}}|p}=3$ : on met la courbe sour la forme (6.1) ; résolvant (6.2) pour $a'_4=1$ et $a'_6=0$, on met la courbe sous la forme type%
% archive sheet number: 28
